% Slide Beamer: GIT -- GitHub -- GitLab cho hoc phan ASPX (.NET Framework 4.8.1)
% File goc (root): nap 8 fragment phan1..phan8 qua \input. Bien dich: pdflatex slide.tex (2 lan).
% Compiler: pdfLaTeX -- TeX Live / MiKTeX / Overleaf; ti le 16:9 cho man hinh rong.
\documentclass[10pt,aspectratio=169]{beamer}
\usetheme{Madrid}
\usepackage[utf8]{inputenc}
\usepackage[T5]{fontenc}
\usepackage[vietnamese]{babel}
\usepackage{lmodern}
\usepackage{listings}
\usepackage{tikz}
\usepackage{booktabs}
\usepackage{hyperref}
\usetikzlibrary{arrows.meta,positioning}

% Cau hinh listings dung chung cho cac frame fragile (vi du lenh bash, chi ASCII).
\lstset{basicstyle=\ttfamily\footnotesize,breaklines=true}
\lstset{keywordstyle=\color{blue},commentstyle=\color{gray},showstringspaces=false}
\lstset{numbers=none,frame=single,tabsize=2,keepspaces=true}

% Metadata: tieu de, tac gia, don vi, ngay trinh bay.
\title[GIT -- GitHub -- GitLab]{GIT -- GitHub -- GitLab \\ cho học phần ASPX (.NET Framework 4.8.1)}
\author{Lê Quang An (0024415453), Lê Tuấn Kiệt (0024417193), Đoàn Minh Trí (0024417136)}
\institute{Học phần ASPX -- Nhóm 3 thành viên}
\date{\today}

% Chan slide gon: chi hien so frame; an bieu tuong dieu huong cho sach.
\setbeamertemplate{footline}[frame number]
\setbeamertemplate{navigation symbols}{}

\begin{document}

% Slide tieu de mo dau bai trinh bay.
\frame{\titlepage}

% Muc luc tong quan 8 phan cua bai trinh bay.
\begin{frame}{Mục lục}
\tableofcontents
\end{frame}

% Phan 1: gioi thieu de tai, muc tieu, pham vi, lien he hoc phan ASPX.
\section{Mở đầu}
\begin{frame}{Giới thiệu nhóm 3 thành viên}
\begin{itemize}
\item Nhóm 3 thành viên, mỗi người phụ trách 1 mảng rõ ràng.
\item Mọi ví dụ gắn với project Web Forms trên .NET Framework 4.8.1.
\item Làm việc theo nhánh \texttt{feature} và review chéo qua pull request.
\item Bảng dưới đây chốt vai trò dùng xuyên suốt bài nói và demo.
\end{itemize}
\begin{center}
\begin{tabular}{lll}
\toprule
Thành viên & Mảng phụ trách & Demo chính \\
\midrule
Lê Quang An (TV1) & Git cục bộ và Internals & Nhánh, merge, conflict Web.config \\
Lê Tuấn Kiệt (TV2) & GitHub và Actions & Pull request, build msbuild \\
Đoàn Minh Trí (TV3) & GitLab và tổng hợp & Merge request, pipeline, so sánh \\
\bottomrule
\end{tabular}
\end{center}
\end{frame}

\begin{frame}{Vì sao đồ án ASPX cần Git?}
\begin{itemize}
\item Gửi ZIP qua Drive hay chat nên dễ đè mất code của nhau.
\item Sửa Web.config rồi mất bản chạy được, không biết cách hoàn tác.
\item Xóa nhầm thư mục hay hỏng máy là mất cả đồ án vì không có lịch sử.
\item Không biết ai sửa dòng nào nên khó truy vết lỗi và review.
\item Máy này chạy được nhưng máy thầy build lỗi mà không tái hiện được.
\end{itemize}
\end{frame}

\begin{frame}{Mục tiêu của đề tài}
\begin{itemize}
\item Nắm vững Git cục bộ: \texttt{commit}, \texttt{branch}, \texttt{merge} và hoàn tác an toàn.
\item Hiểu Git Internals: blob, tree, \texttt{commit} và định danh SHA.
\item Thành thạo GitHub: pull request, review và Actions build ASPX.
\item Thành thạo GitLab: merge request, approval và pipeline Windows runner.
\item Áp dụng quy trình 3 người cho đồ án Web Forms 4.8.1.
\item Có repo demo thật với lịch sử mẫu và CI chạy xanh.
\end{itemize}
\end{frame}

\begin{frame}{Roadmap bài nói}
\begin{enumerate}
\item Phần 1: Mở đầu, mục tiêu và quy ước demo chung.
\item Phần 2: Git từ căn bản đến nâng cao cho project ASPX.
\item Phần 3: GitHub, pull request và Actions build msbuild.
\item Phần 4: GitLab, pipeline và bảng so sánh GitHub với GitLab.
\item Demo live: nhánh, pull request, pipeline xanh và xử lý conflict.
\item Kết luận, bài học nhóm và phần hỏi đáp Q\&A.
\end{enumerate}
\end{frame}

\begin{frame}[fragile]{Quy ước demo chung}
\begin{itemize}
\item Dùng chung một repo demo AspxDemo\_NET481 cho cả nhóm.
\item Ba nhánh \texttt{feature} song song, mỗi người phụ trách một nhánh.
\item Không đẩy thẳng lên \texttt{main}, luôn đi qua pull request có review.
\item Mọi pull request phải chờ CI chạy xanh mới được merge.
\item Commit ngắn gọn theo Conventional Commits để dễ đọc lịch sử.
\end{itemize}
\begin{lstlisting}[language=bash]
git clone https://github.com/MinhTriTM/AspxDemo_NET481.git
git checkout -b feature/login-aspx
git add Login.aspx Login.aspx.cs
git commit -m "feat: add login form"
git push origin feature/login-aspx
\end{lstlisting}
\end{frame}

% Phan 2: cai dat, cau hinh, vung lam viec, commit, nhanh, merge co ban.
\section{Git căn bản}
\begin{frame}{Git là gì? Vì sao dùng Git?}
\begin{itemize}
  \item Git là hệ quản trị phiên bản \textbf{phân tán}, do Linus Torvalds tạo năm 2005.
  \item Mỗi bản sao là một kho đầy đủ lịch sử, làm việc offline rồi đồng bộ sau.
  \item Mọi thay đổi được ghi theo \texttt{commit}, mỗi commit có mã SHA định danh duy nhất.
  \item Gắn với ASPX: quản lý Web Forms, \texttt{Web.config}, biết rõ ai sửa gì, khi nào.
\end{itemize}
\vspace{0.3em}
\begin{tabular}{lll}
\toprule
Tiêu chí & Tập trung (SVN) & Phân tán (Git) \\
\midrule
Nơi lưu lịch sử & Chỉ trên server & Mọi máy đều đủ \\
Mất mạng & Không commit được & Commit bình thường \\
Sập server & Mất lịch sử & Khôi phục từ máy khác \\
\bottomrule
\end{tabular}
\end{frame}

\begin{frame}{Kiến trúc 3 vùng của Git}
\begin{itemize}
  \item \textbf{Working Directory}: các file đang sửa trực tiếp trên đĩa.
  \item \textbf{Staging Area} (Index): vùng đệm chứa các thay đổi chờ ghi nhận.
  \item \textbf{Repository} (thư mục \texttt{.git}): kho lưu toàn bộ lịch sử commit.
  \item \texttt{git add} đưa file từ Working Directory vào Staging Area.
  \item \texttt{git commit} ghi nội dung Staging Area thành một commit trong Repository.
  \item \texttt{git status} cho biết từng file đang nằm ở vùng nào.
  \item Ví dụ ASPX: sửa \texttt{Default.aspx} rồi add và commit để lưu phiên bản.
\end{itemize}
\end{frame}

\begin{frame}[fragile]{Luồng lệnh Git cơ bản}
\begin{itemize}
  \item Cấu hình danh tính một lần duy nhất trên mỗi máy.
  \item Khởi tạo kho mới bằng \texttt{init}, hoặc tải kho có sẵn bằng \texttt{clone}.
  \item Vòng lặp hằng ngày: sửa file, \texttt{add}, \texttt{commit}, xem \texttt{log}.
\end{itemize}
\begin{lstlisting}[language=bash]
git config --global user.name "Nguyen Van A"
git config --global user.email "a@sv.edu.vn"
git init AspxDemo_NET481
cd AspxDemo_NET481
git add Default.aspx Web.config
git commit -m "feat: them trang Default"
git log --oneline
\end{lstlisting}
\end{frame}

\begin{frame}[fragile]{Nhánh và hợp nhất (branch và merge)}
\begin{itemize}
  \item \textbf{Branch} là con trỏ nhẹ tới một commit, dùng để phát triển song song.
  \item Nhánh \texttt{main} luôn ổn định; mỗi tính năng làm trên nhánh \texttt{feature} riêng.
  \item \texttt{merge} gộp nhánh tính năng trở lại nhánh chính sau khi xong việc.
  \item Nhóm 3 người: mỗi TV một nhánh feature, gộp qua Pull Request.
\end{itemize}
\begin{lstlisting}[language=bash]
git switch -c feature/login
# sua file Login.aspx roi commit
git add Login.aspx
git commit -m "feat: them trang dang nhap"
git switch main
git merge feature/login
git branch -d feature/login
\end{lstlisting}
\end{frame}

\begin{frame}{Xung đột (conflict) và cách giải quyết}
\begin{itemize}
  \item Conflict xảy ra khi 2 nhánh cùng sửa một chỗ của một file rồi merge.
  \item Git đánh dấu vùng xung đột bằng \texttt{HEAD}, \texttt{====} và tên nhánh.
\end{itemize}
Quy trình xử lý 4 bước:
\begin{enumerate}
  \item Mở file, đọc kỹ 2 phía và giữ lại nội dung đúng, xóa dấu đánh dấu.
  \item Kiểm tra kỹ, đặc biệt với file \texttt{Web.config} và chuỗi kết nối.
  \item Đánh dấu đã xong bằng lệnh \texttt{git add} trên file vừa sửa.
  \item Hoàn tất bằng \texttt{git commit} (hoặc \texttt{git merge --continue}).
\end{enumerate}
\end{frame}

\begin{frame}[fragile]{Tệp gitignore cho project ASP.NET}
\begin{itemize}
  \item \texttt{.gitignore} liệt kê file và thư mục Git phải bỏ qua.
  \item Project ASPX 4.8.1 phải bỏ qua \texttt{bin}, \texttt{obj}, \texttt{packages} và file user.
  \item Không commit chuỗi kết nối thật và file \texttt{.user} của Visual Studio.
  \item File đã commit rồi thì \texttt{.gitignore} không còn tác dụng với nó.
\end{itemize}
\begin{lstlisting}[language=bash]
# Build output
[Bb]in/
[Oo]bj/
# NuGet packages
packages/
*.nupkg
# Visual Studio user files
*.user
*.suo
.vs/
# Publish output
PublishProfiles/
App_Data/*.mdf
App_Data/*.ldf
\end{lstlisting}
\end{frame}

\begin{frame}{Demo live của TV1: checklist 6 việc}
\begin{itemize}
  \item Demo trên project \texttt{AspxDemo} NET Framework 4.8.1 thật, khoảng 5--7 phút.
\end{itemize}
\begin{enumerate}
  \item \texttt{init} và \texttt{clone} kho demo, cấu hình danh tính và remote.
  \item Sửa \texttt{Default.aspx}, thực hiện \texttt{add}, \texttt{commit}, xem \texttt{log}.
  \item Tạo nhánh \texttt{feature}, commit riêng rồi \texttt{merge} về \texttt{main}.
  \item Tạo conflict cố ý trên \texttt{Web.config} và giải quyết đủ 4 bước.
  \item Demo \texttt{stash}, \texttt{reset} và \texttt{revert} khi sửa nhầm code.
  \item Gắn \texttt{tag} phiên bản v1.0 và đẩy toàn bộ lên GitHub.
\end{enumerate}
\end{frame}


% Phan 3: rebase, stash, reflog, hooks, GPG, GitFlow, Conventional Commits.
\section{Git nâng cao}
\begin{frame}[fragile]{Bên trong Git: blob, tree, commit và SHA}
\begin{itemize}
\item \texttt{blob} lưu nội dung file, \texttt{tree} lưu cây thư mục, \texttt{commit} trỏ tới một \texttt{tree}.
\item Mỗi \texttt{commit} gồm tác giả, thời gian, thông điệp và SHA của \texttt{commit} cha.
\item SHA-1 định danh mỗi đối tượng: đổi 1 byte thì SHA đổi hoàn toàn.
\item Ví dụ ASPX: sửa \texttt{Default.aspx} sẽ tạo \texttt{blob} mới và \texttt{commit} mới.
\item Kiểm tra đối tượng bằng lệnh \texttt{cat-file} như ví dụ dưới đây.
\end{itemize}
\begin{lstlisting}[language=bash]
git cat-file -p HEAD
git cat-file -t HEAD
git hash-object Default.aspx
\end{lstlisting}
\end{frame}

\begin{frame}{Rebase và merge}
\begin{itemize}
\item \texttt{merge} giữ nguyên lịch sử và tạo thêm một \texttt{merge commit}.
\item \texttt{rebase} viết lại lịch sử thành một đường thẳng, log sạch hơn.
\item \texttt{merge} an toàn cho nhánh public vì không đổi SHA cũ.
\item \texttt{rebase} chỉ dùng cho nhánh riêng chưa chia sẻ với ai.
\item Cấm \texttt{rebase} nhánh public: đổi SHA, gây xung đột cho cả nhóm.
\end{itemize}
\begin{tabular}{lll}
\toprule
Tiêu chí & merge & rebase \\
\midrule
Lịch sử & phân nhánh, trung thực & đường thẳng, sạch \\
Merge commit & có & không \\
Xung đột & giải một lần & giải lại từng commit \\
Dùng khi nào & nhánh public & nhánh riêng \\
\bottomrule
\end{tabular}
\end{frame}

\begin{frame}[fragile]{Cherry-pick và bisect}
\begin{itemize}
\item \texttt{cherry-pick} sao chép một \texttt{commit} sang nhánh hiện tại.
\item Dùng khi cần lấy một bản sửa lỗi mà không \texttt{merge} cả nhánh.
\item \texttt{bisect} tìm \texttt{commit} gây lỗi bằng tìm kiếm nhị phân.
\item Kết hợp \texttt{bisect run} với lệnh build để tìm lỗi tự động.
\end{itemize}
\begin{lstlisting}[language=bash]
git cherry-pick a1b2c3d
git cherry-pick --continue
\end{lstlisting}
\begin{lstlisting}[language=bash]
git bisect start
git bisect bad
git bisect good v1.0
git bisect run msbuild AspxDemo.sln
git bisect reset
\end{lstlisting}
\end{frame}

\begin{frame}[fragile]{Hooks và LFS}
\begin{itemize}
\item \texttt{Hook} là script tự chạy khi \texttt{commit}, \texttt{push} hay \texttt{merge}.
\item Ví dụ: \texttt{pre-commit} kiểm tra build, \texttt{commit-msg} kiểm tra định dạng.
\item \texttt{LFS} lưu file lớn như DLL hay ảnh ngoài repo, giữ repo nhẹ.
\item Trong repo chỉ là con trỏ text, nội dung thật nằm trên server LFS.
\end{itemize}
\begin{lstlisting}[language=bash]
ls .git/hooks
git lfs track "*.dll"
git lfs ls-files
\end{lstlisting}
\end{frame}

\begin{frame}{Quy trình nhánh: GitFlow, GitHub Flow, Trunk-Based}
\begin{itemize}
\item GitFlow nhiều nhánh, hợp với phát hành theo phiên bản lớn.
\item GitHub Flow đơn giản: một nhánh \texttt{main} cộng feature và PR.
\item Trunk-Based commit nhỏ liên tục lên \texttt{main}, đòi hỏi CI mạnh.
\item Nhóm 3 người làm ASPX 4.8.1 nên chọn GitHub Flow.
\item Mọi thay đổi đi qua PR với ít nhất một người review.
\end{itemize}
\begin{tabular}{lll}
\toprule
Quy trình & Nhánh chính & Phù hợp \\
\midrule
GitFlow & main, develop, feature & phát hành theo bản \\
GitHub Flow & main, feature & nhóm nhỏ, web \\
Trunk-Based & main & CI mạnh, nhịp nhanh \\
\bottomrule
\end{tabular}
\end{frame}

\begin{frame}{Conventional Commits}
\begin{itemize}
\item Định dạng chung: \texttt{type(scope): mô tả ngắn gọn}.
\item \texttt{feat(auth): thêm trang đăng nhập Web Forms}
\item \texttt{fix(config): sửa connection string Web config}
\item \texttt{docs(slide): bổ sung hình git flow}
\item \texttt{test(build): thêm test msbuild pipeline}
\item Lợi ích: log dễ đọc, dễ lọc và tự sinh changelog.
\end{itemize}
\end{frame}


% Phan 4: repo, Issues, Fork, Pull Request, review, branch protection.
\section{GitHub}
\begin{frame}{GitHub là gì?}
\begin{itemize}
\item Nền tảng hosting Git trực tuyến lớn nhất cho mã nguồn mở và doanh nghiệp.
\item Quản lý repository tập trung với giao diện web và phân quyền theo vai trò.
\item Cộng tác qua pull request, fork và review code ngay trên trình duyệt.
\item Quản lý công việc bằng Issues, Projects và Wiki cho đồ án nhóm.
\item Tự động build, test và triển khai với Actions pipelines tích hợp sẵn.
\item Bảo mật và phân phối: Dependabot, secret scanning, Pages và Packages.
\end{itemize}
\end{frame}

\begin{frame}{Issues, Projects và Wiki}
\begin{itemize}
\item \textbf{Issues}: thẻ việc có người phụ trách, nhãn và milestone rõ ràng.
\item \textbf{Projects}: bảng Kanban kéo thả thẻ việc qua Todo, Doing và Done.
\item Ví dụ đồ án: Issue \texttt{fix-login} mô tả lỗi form \texttt{Login.aspx} kèm ảnh chụp.
\item Ví dụ đồ án: thẻ Project kiểm tra tích hợp file \texttt{Web.config} mẫu trước khi merge.
\item Ví dụ đồ án: Wiki ghi quy trình build \texttt{msbuild} và cấu hình IIS cho 4.8.1.
\end{itemize}
\end{frame}

\begin{frame}{Fork và branch}
\begin{itemize}
\item \textbf{Fork}: bản sao repository độc lập thuộc tài khoản khác.
\item \textbf{Branch}: nhánh nhẹ trong cùng repository để làm tính năng.
\item Đồ án nhóm 3 người ưu tiên branch vì review nhanh và chung pipeline.
\item Fork phù hợp nhận đóng góp từ cộng đồng mã nguồn mở bên ngoài.
\end{itemize}
\smallskip
{\small
\centering
\begin{tabular}{lll}
\toprule
Tiêu chí & Fork & Branch \\
\midrule
Đối tượng & Người ngoài dự án & Thành viên nhóm \\
Không gian & Repository mới & Cùng repository \\
Đồng bộ & Pull từ upstream & Merge nội bộ \\
Quyền ghi & Không cần quyền gốc & Cần quyền ghi nhánh \\
\bottomrule
\end{tabular}
}
\end{frame}

\begin{frame}{Vòng đời Pull Request}
\begin{enumerate}
\item Tạo nhánh tính năng và đẩy commit lên GitHub.
\item Mở pull request về \texttt{main} kèm mô tả và ảnh giao diện.
\item Review code, thảo luận và đẩy thêm commit sửa lỗi.
\item Chờ pipeline Actions chạy xanh và đủ lượt approve.
\item Merge pull request rồi xóa nhánh tính năng đã hoàn thành.
\end{enumerate}
\end{frame}

\begin{frame}[fragile]{Review code và protected branch}
\begin{itemize}
\item Mọi đổi mới vào \texttt{main} phải qua pull request có ít nhất một lượt review.
\item Reviewer kiểm tra logic, giao diện và file \texttt{Web.config} mẫu an toàn.
\item Bật required status checks để chặn merge khi pipeline Actions báo đỏ.
\item Chặn push trực tiếp và force push lên nhánh \texttt{main} của đồ án.
\end{itemize}
\begin{lstlisting}[language=bash]
# Rule: protect main branch
# - require 1 approval before merge
# - require check ci-aspx-481 to pass
# - block force push to main
\end{lstlisting}
\end{frame}

\begin{frame}[fragile]{CODEOWNERS tự động gán reviewer}
\begin{itemize}
\item File \texttt{CODEOWNERS} đặt ở gốc hoặc thư mục \texttt{docs} của repository.
\item Mỗi dòng gán một mẫu đường dẫn cho người hoặc nhóm review phụ trách.
\item Pull request chạm file nào thì owner tương ứng được yêu cầu review.
\item Kết hợp protected branch để bắt buộc approve của owner trước khi merge.
\end{itemize}
\begin{lstlisting}[language=bash]
*.aspx @nhom3/tv-frontend
Web.config @nhom3/tv-backend
.github/workflows/ @nhom3/tv-devops
\end{lstlisting}
\end{frame}


% Phan 5: pipeline CI build msbuild + test cho project ASPX 4.8.1.
\section{CI với GitHub Actions}
% Fragment: Phan 5 - GitHub Actions cho ASPX 4.8.1 (TV2)
% File nay chi chua cac frame Beamer, duoc \input vao slide.tex.

\begin{frame}{GitHub Actions: workflow, job, step, runner}
\begin{itemize}
\item \textbf{Workflow}: tệp YAML trong \texttt{.github/workflows}, tự chạy khi \texttt{push}, mở \texttt{PR} hoặc theo lịch.
\item \textbf{Job}: nhóm các step chạy trên cùng một runner; các job chạy song song, dùng \texttt{needs} để xếp nối tiếp.
\item \textbf{Step}: từng bước nhỏ trong job, dùng \texttt{uses} gọi action có sẵn hoặc \texttt{run} chạy lệnh shell.
\item \textbf{Runner}: máy ảo thực thi job; ASPX 4.8.1 bắt buộc \texttt{windows-latest} vì cần \texttt{msbuild} bản full.
\end{itemize}
\end{frame}

\begin{frame}[fragile]{Pipeline CI cho ASPX 4.8.1}
\begin{itemize}
\item Mọi \texttt{push} và \texttt{PR} đều kích hoạt build và test tự động.
\item \texttt{msbuild} biên dịch \texttt{AspxDemo\_NET481.sln}, \texttt{vstest} chạy test, artifact lưu gói web.
\end{itemize}
\begin{lstlisting}[basicstyle=\ttfamily\scriptsize]
name: CI ASPX 4.8.1
on: [push, pull_request]
jobs:
  build:
    runs-on: windows-latest
    steps:
      - uses: actions/checkout@v4
      - name: Setup MSBuild
        uses: microsoft/setup-msbuild@v2
      - name: Setup NuGet
        uses: NuGet/setup-nuget@v2
      - name: Restore packages
        run: nuget restore AspxDemo_NET481.sln
      - name: Build solution
        run: msbuild AspxDemo_NET481.sln /p:Config=Release
      - name: Run tests
        run: vstest.console.exe Tests.dll
      - name: Upload artifact
        uses: actions/upload-artifact@v4
        with: {name: webapp, path: bin/}
\end{lstlisting}
\end{frame}

\begin{frame}{Vì sao ASPX 4.8.1 bắt buộc Windows runner}
\begin{itemize}
\item .NET Framework 4.8.1 và \texttt{msbuild} bản full chỉ có trên Windows.
\item Lệnh \texttt{dotnet} trên Ubuntu không biên dịch được Web Forms cũ.
\item Chốt: mọi job build ASPX đều đặt \texttt{windows-latest}.
\end{itemize}
\begin{center}
\small
\begin{tabular}{lcc}
\toprule
Tiêu chí & Windows & Ubuntu \\
\midrule
MSBuild + Framework 4.8.1 & Có sẵn & Không có \\
Build Web Forms cũ & Được & Không được \\
IIS, \texttt{vstest} & Hỗ trợ & Không hỗ trợ \\
Tốc độ, chi phí & Chậm, đắt hơn & Nhanh, rẻ \\
\bottomrule
\end{tabular}
\end{center}
\end{frame}

\begin{frame}[fragile]{Secrets và Environments}
\begin{itemize}
\item \textbf{Secret}: biến mã hóa trong Settings, không bao giờ in ra log (bị che thành \texttt{***}).
\item \textbf{Environment}: nhóm secret theo \texttt{dev}, \texttt{staging}, \texttt{prod}, kèm người duyệt deploy.
\item Không commit connection string vào \texttt{Web.config}; secret được thay đúng lúc deploy.
\item Cách đọc secret trong workflow và lệnh build:
\end{itemize}
\begin{lstlisting}[language=bash, basicstyle=\ttfamily\scriptsize]
echo "Checking ${{ secrets.DB_SERVER }}"
nuget restore AspxDemo_NET481.sln
msbuild AspxDemo_NET481.sln /p:Conn="${{ secrets.CONN_STR }}"
\end{lstlisting}
\end{frame}

\begin{frame}{Demo TV2: PR và pipeline xanh}
\begin{enumerate}
\item Tạo nhánh \texttt{feature/github-actions}, thêm file workflow, push lên GitHub.
\item Mở \texttt{PR} sang \texttt{main}, gắn TV1 hoặc TV3 làm reviewer.
\item Reviewer review code rồi Approve; chờ Actions chạy xong báo tick xanh.
\item Merge \texttt{PR} rồi xóa nhánh; kiểm tra \texttt{main} vẫn xanh.
\item Chiếu pipeline xanh và gói artifact lên màn hình lớp.
\end{enumerate}
\end{frame}

\begin{frame}{Bảo mật: secret scanning, Dependabot, CodeQL}
\begin{itemize}
\item \textbf{Secret scanning}: quét token và key lộ trong code, chặn push chứa secret mới.
\item \textbf{Dependabot}: tự mở \texttt{PR} nâng cấp gói NuGet dính lỗi CVE.
\item \textbf{CodeQL}: phân tích tĩnh, tìm SQL injection và XSS trong code C\#.
\item Bật cả ba trong mục Security; job thất bại khi phát hiện lỗi nghiêm trọng.
\end{itemize}
\end{frame}


% Phan 6: Groups, Merge Request, Approval, GitLab CI tren Windows runner.
\section{GitLab}
\begin{frame}{GitLab là gì? SaaS vs self-hosted}
\begin{itemize}
\item Nền tảng DevOps all-in-one trên Git: repo, Merge Request, Issues, CI/CD tích hợp sẵn.
\item Một vòng đời khép kín: plan, code, review, test, release, monitor trong cùng project.
\item Dùng ngay bản SaaS tại \texttt{gitlab.com} hoặc tự vận hành bản self-hosted.
\end{itemize}
\begin{center}
\small
\begin{tabular}{lll}
\toprule
Tiêu chí & SaaS & Self-hosted \\
\midrule
Hạ tầng & GitLab quản lý & Nhóm tự cài đặt \\
Chi phí \& bảo trì & Trả theo gói, không lo máy chủ & Miễn phí bản Community, tốn công admin \\
Phù hợp & Đồ án môn học, demo nhanh & Doanh nghiệp cần kiểm soát dữ liệu \\
\bottomrule
\end{tabular}
\end{center}
\end{frame}

\begin{frame}{Groups, Projects, Issues, Boards}
\begin{itemize}
\item Groups gom nhiều project theo nhóm/lớp: phân quyền member, chia sẻ runner chung.
\item Projects chứa repo Git, Merge Request, pipeline, wiki; gắn ASPX qua \texttt{main} bảo vệ.
\item Issues theo dõi task/bug với label, milestone, assignee; đóng bằng \texttt{Closes} khi merge.
\item Boards kiểu Kanban kéo thả Issues qua các cột To Do, Doing, Done theo milestone.
\end{itemize}
\end{frame}

\begin{frame}{Merge Request và Approval rules}
\begin{enumerate}
\item Tạo nhánh \texttt{feature} từ \texttt{main}, push và mở Merge Request kèm mô tả rõ ràng.
\item Hệ thống chạy pipeline kiểm tra build và test tự động trên nhánh nguồn.
\item Reviewer nhận xét từng dòng, yêu cầu chỉnh sửa cho đến khi đạt yêu cầu.
\item Đủ lượt approve theo rule mới được phép merge, đảm bảo ít nhất 2 người đồng ý.
\item Nhấn merge, xóa nhánh nguồn, Issues liên quan tự đóng theo từ khóa.
\end{enumerate}
\vspace{0.3em}
Ví dụ rule: nhánh \texttt{main} yêu cầu 2 approve từ maintainer và pipeline phải xanh 100\%.
\end{frame}

\begin{frame}[fragile]{Tệp gitlab-ci YAML cho ASPX 4.8.1}
\begin{itemize}
\item Khai báo \texttt{stages}, \texttt{variables}, job theo từng giai đoạn build, test, deploy.
\end{itemize}
\begin{lstlisting}[language=bash]
stages:
  - build
  - test
  - deploy
variables:
  SOLUTION: AspxDemo.sln
  CONFIG: Release
build:aspx:
  stage: build
  tags: [windows-481]
  script:
    - nuget restore %SOLUTION%
    - msbuild %SOLUTION% /p:Config=%CONFIG%
  artifacts:
    paths: [bin/Release/]
    expire_in: 1 week
test:aspx:
  stage: test
  tags: [windows-481]
  script: [vstest.console.exe bin/Release/*.Tests.dll]
\end{lstlisting}
\end{frame}

\begin{frame}{Runners và tags Windows}
\begin{itemize}
\item Runner là agent chạy job: shell runner trên Windows Server cài Visual Studio và \texttt{msbuild}.
\item Tags điều phối job đúng máy: job ASPX 4.8.1 gắn tag \texttt{windows-481}, runner cùng tag mới nhận.
\item Runner nhóm dùng chung cho 3 thành viên; bật cache NuGet để build lần sau nhanh hơn.
\end{itemize}
\end{frame}

\begin{frame}{Environments và Registry}
\begin{itemize}
\item Environments quản lý staging và production, hỗ trợ deploy thủ công và rollback một nút bấm.
\item Review Apps tự tạo môi trường xem trước cho mỗi Merge Request, xóa khi merge xong.
\item Container Registry lưu image Docker và gói NuGet nội bộ, dùng token nhóm để pull an toàn.
\end{itemize}
\end{frame}

\begin{frame}{Demo TV3: MR và pipeline}
\begin{itemize}
\item[$\square$] Mirror repo ASPX sang GitLab, bảo vệ nhánh \texttt{main} và bật Approval 2 người.
\item[$\square$] Mở Merge Request từ \texttt{feature} kèm mô tả, label và người review cụ thể.
\item[$\square$] Theo dõi pipeline build \texttt{msbuild} trên runner Windows đến khi báo xanh.
\item[$\square$] Thu đủ approve, merge vào \texttt{main} và kiểm tra deploy staging thành công.
\item[$\square$] Chụp pipeline xanh và bảng so sánh đưa vào báo cáo Chương 5 và Chương 6.
\end{itemize}
\end{frame}


% Phan 7: bang so sanh GitHub -- GitLab va quy trinh nhom 3 nguoi.
\section{So sánh và quy trình nhóm}
\begin{frame}{Git, GitHub, GitLab khác nhau thế nào?}
\begin{itemize}
\item Git là công cụ quản lý phiên bản phân tán, chạy trên máy cá nhân.
\item GitHub và GitLab là máy chủ chứa repo, cộng thêm cộng tác và CI/CD.
\item Một project dùng Git ở local, đẩy lên GitHub hoặc GitLab để chia sẻ.
\item Hiểu đúng: Git là nền tảng bắt buộc, GitHub và GitLab là lựa chọn bổ sung.
\end{itemize}
{\small
\begin{tabular}{@{}p{2.0cm}p{3.4cm}p{4.2cm}@{}}
\toprule
Khía cạnh & Git & GitHub, GitLab \\
\midrule
Bản chất & Tool dòng lệnh & Dịch vụ hosting và cộng tác \\
Nơi chạy & Máy cá nhân & Cloud hoặc self-host \\
Lịch sử & Đầy đủ trong .git & Nhân bản cộng thêm PR, MR, CI \\
Ví dụ & \texttt{git log} tại local & Mở PR trên web để review \\
\bottomrule
\end{tabular}}
\end{frame}

\begin{frame}{GitHub và GitLab: 8 tiêu chí cho đồ án}
\begin{itemize}
\item Bảng dưới so sánh 8 điểm ảnh hưởng trực tiếp tới đồ án ASPX.
\item Nhóm môn học ưu tiên: miễn phí, CI dễ học, review nhanh.
\item Cả hai đều hỗ trợ repo private miễn phí cho sinh viên.
\item Chốt nhanh: chọn một nền tảng chính, mirror sang nền tảng còn lại.
\end{itemize}
{\footnotesize
\renewcommand{\arraystretch}{0.9}
\begin{tabular}{@{}p{2.6cm}p{3.3cm}p{3.7cm}@{}}
\toprule
Tiêu chí & GitHub & GitLab \\
\midrule
Giá sinh viên & Free và Education pack & Free và Ultimate giáo dục \\
CI cho ASPX 4.8.1 & Actions, runner Windows & Pipeline, runner Windows \\
Review code & PR và rulesets & MR và approval rules \\
Quản lý việc & Issues, Projects & Issues, Boards, Epics \\
Bảo mật & Dependabot, CodeQL & SAST, secret detection \\
Self-host & Không, chỉ cloud & Có, bản community \\
Tài liệu mẫu & Rất nhiều ví dụ & Đầy đủ, gắn DevOps \\
Phù hợp ASPX & Tốt, Actions dễ học & Tốt, pipeline mạnh \\
\bottomrule
\end{tabular}}
\end{frame}

\begin{frame}[fragile]{Nên chọn nền tảng nào cho đồ án?}
\begin{itemize}
\item Khuyến nghị 1: chọn GitHub làm chính nếu nhóm mới học CI lần đầu.
\item Khuyến nghị 2: chọn GitLab làm chính nếu muốn pipeline nhiều stage và review chặt.
\item Khuyến nghị 3: dùng chính một nơi, mirror sang nơi còn lại để dự phòng.
\item Mọi ví dụ ASPX 4.8.1 đều build bằng runner Windows và \texttt{msbuild}.
\item Thầy cô chấm nhanh nhờ link PR hoặc MR cộng với pipeline xanh.
\end{itemize}
\begin{lstlisting}[language=bash]
# mirror GitHub -> GitLab (backup)
git remote add gitlab <gitlab-url>
git push gitlab main --tags
\end{lstlisting}
\end{frame}

\begin{frame}[fragile]{Quy trình 3 người: main, develop, feature}
\begin{itemize}
\item Luật 1: cấm push thẳng \texttt{main} và \texttt{develop}, mọi đổi phải qua PR hoặc MR.
\item Luật 2: mỗi PR cần một người review và pipeline xanh mới được merge.
\item Luật 3: đặt tên nhánh \texttt{feature/ten-viec}, commit theo Conventional Commits.
\item TV1 giữ \texttt{develop} ổn định, TV2 và TV3 làm trên nhánh feature riêng.
\item Conflict file \texttt{Web.config}: người tạo PR tự merge bản mới nhất rồi sửa.
\end{itemize}
\begin{lstlisting}[language=bash]
# so do nhanh: main <- develop <- feature
main    :  o--------o--------o  (release)
develop :  o---o---o----o---o  (integration)
feature :   \o/      \o/  (1 branch / person)
\end{lstlisting}
\end{frame}

\begin{frame}[fragile]{Review checklist và quản lý secret}
\begin{columns}[t]
\column{0.48\textwidth}
Checklist review:
\begin{itemize}
\item Code build sạch, không warning mới, test chính pass.
\item Không còn marker conflict, đã tự review diff trước khi nhờ review.
\item Cập nhật ghi chú ngắn để thầy cô hiểu đổi gì, vì sao.
\end{itemize}
\column{0.48\textwidth}
Quản lý secret:
\begin{itemize}
\item Không commit connection string thật, mật khẩu, file publish.
\item Dùng biến môi trường và secret của CI cho \texttt{Web.config} khi build.
\item Lỡ lộ secret: thu hồi ngay, xóa lịch sử, báo thầy cô.
\end{itemize}
\end{columns}
\begin{lstlisting}[language=bash]
# remove leaked config from history
git rm --cached Web.config
git commit -m "fix: stop tracking local config"
\end{lstlisting}
\end{frame}

\begin{frame}[fragile]{Bài tập về nhà}
\begin{itemize}
\item Bài 1 (Git local): tạo nhánh \texttt{feature/them-trang}, commit 3 lần, merge vào \texttt{develop}.
\item Bài 2 (PR và CI): mở PR, nhờ bạn review, pipeline build \texttt{msbuild} phải xanh.
\item Bài 3 (Secret): chuyển connection string vào secret của CI, chứng minh lịch sử sạch.
\item Nộp link PR hoặc MR kèm ảnh pipeline xanh trước buổi sau.
\item Chấm điểm: đúng quy trình 40, CI xanh 30, review chất 30.
\end{itemize}
\begin{lstlisting}[language=bash]
# starter for homework
git checkout -b feature/them-trang
git log --oneline -5
git push -u origin feature/them-trang
\end{lstlisting}
\end{frame}

% Phan 8: tong ket bai hoc, han che, huong phat trien va hoi dap.
\section{Kết luận và hỏi đáp}
\begin{frame}{Kết luận: 5 takeaway chính}
\begin{itemize}
\item Git lưu vết mọi thay đổi: \texttt{commit} nhỏ, message rõ, lịch sử sạch.
\item Nhánh và \texttt{merge} giữ \texttt{main} ổn định, conflict xử lý sớm.
\item GitHub và GitLab biến Git thành cộng tác: PR, MR, review, Issues.
\item CI với \texttt{msbuild} và test tự động: pipeline xanh mới được merge.
\item Quy trình 3 người áp dụng được ngay cho đồ án ASPX 4.8.1.
\end{itemize}
\end{frame}

\begin{frame}{Mỗi thành viên một câu takeaway}
\begin{itemize}
\item Mỗi thành viên đúc kết một câu sau khi làm đồ án này.
\item TV1 -- Git local: hiểu 3 vùng và \texttt{commit} thì không còn sợ mất code.
\item TV2 -- GitHub: PR và Actions giúp review kỹ và build ASPX tự động.
\item TV3 -- GitLab: MR và pipeline giúp nhóm nhỏ làm việc như nhóm chuyên nghiệp.
\item Điểm chung: kỷ luật \texttt{commit} và review quan trọng hơn công cụ.
\end{itemize}
\end{frame}

\begin{frame}{Đóng góp từng thành viên}
\begin{itemize}
\item Phân công theo mảng: Git, GitHub, GitLab và tổng hợp.
\item Mỗi người viết báo cáo, slide và demo live của mảng mình.
\item Review chéo qua PR và MR, không đẩy thẳng lên \texttt{main}.
\item Bảng dưới đây tóm tắt mảng và demo của từng thành viên.
\end{itemize}
{\footnotesize
\begin{tabular}{lll}
\toprule
Thành viên & Mảng phụ trách & Demo live \\
\midrule
TV1 & Git cơ bản và nâng cao & init, branch, merge, conflict Web.config \\
TV2 & GitHub và Actions & PR, review, Actions build msbuild \\
TV3 & GitLab và so sánh & MR, approval, pipeline, bảng so sánh \\
\bottomrule
\end{tabular}
}
\end{frame}

\begin{frame}[fragile]{Repo demo và pipeline xanh}
\begin{itemize}
\item Repo GitHub: \url{https://github.com/nhom-aspx/AspxDemo-NET481}
\item Repo GitLab: \url{https://gitlab.com/nhom-aspx/AspxDemo-NET481}
\item Pipeline GitHub Actions và GitLab CI đều xanh trước khi merge.
\item Clone repo về máy và xem lịch sử như sau:
\end{itemize}
\begin{lstlisting}[language=bash]
git clone https://github.com/nhom-aspx/AspxDemo-NET481.git
cd AspxDemo-NET481
git log --oneline -5
git branch -a
\end{lstlisting}
\end{frame}

\begin{frame}{Cảm ơn thầy cô và các bạn!}
\begin{center}
{\Large Cảm ơn thầy cô và các bạn đã lắng nghe!}\\
\vspace{0.4cm}
{\large Nhóm 3 -- GIT, GitHub, GitLab cho ASPX 4.8.1}\\
\vspace{0.4cm}
{\large Mời thầy cô và các bạn đặt câu hỏi}\\
\vspace{0.4cm}
Git $\rightarrow$ TV1 \quad GitHub $\rightarrow$ TV2 \quad GitLab $\rightarrow$ TV3
\end{center}
\end{frame}


\end{document}
